\markright{\mititulo \hspace{5cm} \thepage}
\pagestyle{fancy}

\section*{Resumen}
\addcontentsline{toc}{section}{Resumen}

\noindent Se replica el experimento de Eratóstenes para medir el radio terrestre
usando un gnomon en la plazoleta central de la Universidad de Antioquia,
Medellín, Colombia ($6{,}267^\circ$N, $75{,}569^\circ$O).
Se realizaron dos sesiones de medición: el 19~de~marzo de~2026 (un día
antes del equinoccio vernal) y el 20~de~marzo de~2026 (equinoccio vernal).
En cada sesión se registró la longitud de la sombra proyectada por el
gnomon a lo largo del mediodía solar; la longitud mínima de la sombra
se usó para calcular la distancia cenital, la cual, en la condición de
equinoccio, es igual a la latitud geográfica del observador.
Combinando ese ángulo con la distancia meridional ecuador--Medellín
(calculada sobre el elipsoide WGS84), se obtuvo un radio terrestre de
$R_{20} = (6\,310 \pm 182)$~km en el equinoccio (error relativo del
0{,}96\,\%) y $R_{19} = (6\,215 \pm 229)$~km el día~19 (2{,}44\,\%).
Ambos resultados contienen el valor de referencia de la IAU
($R = 6\,371$~km) dentro de su intervalo de incertidumbre,
demostrando que es posible reproducir el resultado de Eratóstenes
con materiales básicos y con una precisión inferior al~1\,\%.

\vspace{1cm}\textit{Palabras clave:} Eratóstenes, gnomon, radio terrestre, equinoccio, incertidumbre experimental

\newpage

\section*{Abstract}
\addcontentsline{toc}{section}{Abstract}

\noindent The experiment of Eratosthenes is reproduced to measure Earth's radius
using a gnomon at the central plaza of the Universidad de Antioquia,
Medellín, Colombia ($6.267^\circ$N, $75.569^\circ$W).
Two measurement sessions were carried out: 19~March~2026 (one day before
the vernal equinox) and 20~March~2026 (vernal equinox).
The shadow length cast by the gnomon was recorded around solar noon;
its minimum value was used to compute the zenith distance, which at the
equinox equals the geographic latitude of the observer.
Combining that angle with the meridional distance from the equator to
Medellín (computed on the WGS84 ellipsoid), an Earth radius of
$R_{20} = (6{,}310 \pm 182)$~km was obtained at the equinox
(relative error 0.96\,\%) and $R_{19} = (6{,}215 \pm 229)$~km on
day~19 (2.44\,\%). Both results include the IAU reference value
($R = 6{,}371$~km) within their uncertainty interval, demonstrating
that Eratosthenes' result can be reproduced with basic materials
and a precision below~1\,\%.

\vspace{1cm}\textit{Keywords:} Eratosthenes, gnomon, Earth radius, equinox, experimental uncertainty

\newpage
